%%%BLOCK id=033-01 ch=03 sec=板块模型·轻板模型 kind=problem alt=none cont=none y=0.05-0.37
\textbf{轻板模型}\quad $M_{\text{板}}=0$\quad $F_{\text{合}}=m_{\text{板}}a$\quad $F_{\text{合}}=0\to f_1=f_2$\quad 质小物体\ $M$ 动不动取决于外力大小% 小字“质小物体”写在“M动不动取决于外力大小”上方；原稿此行向右上倾斜

\begin{problem}
%FIG p033_a 0.08 0.10 0.26 0.175
\notefig[0.3]{figures/p033_a.png}
% 图内标注：A(1kg)、B(2kg)、F；A、B 并排放在一块长板上，标 μ_1=0.2(左端)、F_A、2N、一小箭头(←)、B、4N(右端)，板下标 μ_2=0；图上有几处白色涂改斑块
\begin{solution}
$f_{B\max}>f_{A\max}=2\unit{N}$\quad B 与板相对静止\quad B 与板等效成一个有质板\ $m=2\unit{kg}$

$\mu m_Ag\ge m_Ba$

$F=(m+2m)a\le3\unit{N}$\quad B 最多受\ $f=2\unit{N}$

\noindent A、\ $F=1\unit{N}$ 时\ 不动，均静止\ $\checkmark$\\
B、\ $F=4\unit{N}$ 时\ $f_B=2\unit{N}$\ $\checkmark$\\
C、\ $F=6\unit{N}$ 时\ $a_B=2\unit{m/s^2}$\ $\checkmark$\\
D、\ $F=1.5\unit{N}$ 时\ $f_A=1\unit{N}$\ $\checkmark$
\end{solution}
\end{problem}
%%%END
%%%BLOCK id=033-02 ch=03 sec=斜面·丝带模型 kind=problem alt=none cont=none y=0.39-0.70
\begin{problem}
$M>m$

%FIG p033_b 0.055 0.43 0.265 0.50
\notefig[0.35]{figures/p033_b.png}
% 图内标注：屋脊形斜面上铺一条丝带，左侧放 M(较大)，右侧放 m(较小)，两侧斜面底角均标 θ，右下标“光滑丝带”(“光滑”二字不太确定)
\begin{solution}
\noindent\underline{$M$ 和 $m$ \unclear{与绳带}必相对静止}% 标题带下划线，“和”字后的 m 及“与绳带必”几字辨认不清

\[
M\left\{\begin{array}{@{}l@{}}
\mu>\tan\theta\ \ \text{$M$ 带 $m$ 加速，都相对丝带静止}\\[6pt]
\mu<\tan\theta\ \ \text{$M$ 向左下加速，$m$ 向右下加速，$m$ 与丝带相对滑动}
\end{array}\right.
\]

\noindent 使 $M,m$ 与绳带不发生相对滑动，应满足
\[
\mu mg\cos\theta\ge f
\]
\noindent\hspace*{10em}$\nwarrow$ 共加速时 $f$

\noindent $\mu=\tan\theta$，$m$ 不动，$M$ 下滑\\
\hspace*{10em}与丝带一起
\end{solution}
\end{problem}
%%%END
%%%BLOCK id=033-03 ch=03 sec=连接体·叠放三物体与滑轮 kind=method alt=none cont=none y=0.77-0.96
%FIG p033_c 0.40 0.775 0.665 0.88
\notefig[0.4]{figures/p033_c.png}
% 图内标注：三个物体叠放，1 在 2 上，2 在 3 上；3 受 F 向左；轻绳绕过右端圆柱(滑轮)，上股连 2、下股连 3，两处均标 T；2、3 间摩擦标 f_2，3 与地面间摩擦标 f_3
\noindent 1\ 2 共加速，3 加速
\[
F-f_3-f_2-2T=m_3a
\]% CHECK: 图中作用在 3 上的只有一股绳(T)，此处“2T”存疑，照原样转录
\[
T-f_2=(m_1+m_2)a.
\]
%%%END
%%%BLOCK id=034-01 ch=03 sec=板块模型·有质板块·平地 kind=method alt=01 cont=none y=0.05-0.23
\textbf{有质板块、}

\textbf{平地}

\noindent 地滑

%FIG p034_a 0.088 0.15 0.215 0.222
\notefig[0.28]{figures/p034_a.png}
% 图内标注：B 在 A 上，B 标“2”、向右箭头 V_0，接触面标 μ_1，A 板，板下标“1”
%FIG p034_b 0.24 0.123 0.425 0.215
\notefig[0.4]{figures/p034_b.png}
% 图内标注：v-t 图，纵轴 V、V_0，下降线标“2”(斜线阴影)，上升线标“1”，共速后水平线标“1,2”，横轴 t
\[
\left\{\begin{aligned}
a_1&=\frac{\mu_1m_2g}{m_1}\\
a_2&=\mu_1g
\end{aligned}\right.
\qquad
\begin{aligned}
V_B&=V_0-\mu_1gt\\
V_1&=a_1t
\end{aligned}
\qquad
\begin{aligned}
S_2&=V_0t-\frac12a_2t^2\\
S_1&=\frac12a_1t^2
\end{aligned}
\]% CHECK: V_B 的下标写得像“B”(其余处用 V_1、V_2)
\[
L=S_{\text{相}}=S_A-S_B
\]% CHECK: 此式在原页位于右上方；照原样转录，按图中 A 为板、B 为块，相对位移似应为 S_B-S_A
%%%END
%%%BLOCK id=034-02 ch=03 sec=板块模型·有质板块·初速类 kind=method alt=01 cont=none y=0.23-0.45
\textbf{初速类}

\noindent 地糙 $\mu_2$\quad ① 底滑面糙\ $\mu_1>\mu_2$

\noindent 共速后一起减速

%FIG p034_c 0.185 0.3235 0.365 0.378
\notefig[0.3]{figures/p034_c.png}
% 图内标注：板(标“1”)上右端放物块(标“2”)，板上表面标 μ_1，下表面标 μ_2，板向右箭头 V_0
%FIG p034_d 0.405 0.275 0.60 0.378
\notefig[0.33]{figures/p034_d.png}
% 图内标注：v-t 图，纵轴 V、V_0，从 V_0 下降线标“1”，从原点上升线标“2”，交点后共同下降线标“1.2”，横轴 t
\begin{align*}
a_1&=\frac{\mu_1m_2g+\mu_2(m_1+m_2)g}{m_1} & V_1&=V_0-a_1t\\
a_2&=\mu_1g & V_2&=a_2t\\
S_1&=V_0t-\frac12a_1t^2 & S_2&=\frac12a_2t^2
\end{align*}
\[
\left\{\begin{aligned}
V_{\text{共}}&=a_2t=V_0-a_1t\\
a_{\text{共}}&=\mu_2g\\
X_{\text{共}}&=\frac{V_{\text{共}}^2}{2a_{\text{共}}}
\end{aligned}\right.
\]
%%%END
%%%BLOCK id=034-03 ch=03 sec=板块模型·有质板块·初速类 kind=method alt=01 cont=none y=0.45-0.68
\noindent ② 底糙面滑\ $\mu_1<\mu_2$

\noindent 共速后分别减速

%FIG p034_e 0.175 0.511 0.365 0.557
\notefig[0.3]{figures/p034_e.png}
% 图内标注：板(标“A”“1”)上右端放物块(标“B”“2”)，板上表面标 μ_1，下表面标 μ_2，板向右箭头 V_0
\begin{align*}
a_1&=\frac{\mu_1m_2g+\mu_2(m_1+m_2)g}{m_1} & V_1&=V_0-a_1t & S_1&=V_0t-\frac12a_1t^2\\
a_2&=\mu_1g & V_2&=a_2t & S_2&=\frac12a_2t^2\\
V_{\text{共}}&=a_2t=V_0-a_1t
\end{align*}
%FIG p034_f 0.11 0.568 0.31 0.675
\notefig[0.33]{figures/p034_f.png}
% 图内标注：v-t 图，纵轴 V，A 的下降线(红笔描过)标“1 A”，上升线标“B2”，交点后两条下降线：较缓的标“2 B”、较陡的(红笔)标“A”，其间阴影，横轴 t
\noindent 分别减速

\noindent 对2：$\mu_1m_2g=m_2a_2'$\quad $a_2'=\mu_1g$\\
对1：$\mu_1m_2g-\mu_2(m_1+m_2)g=m_1a_1'$\quad $a_1'=\cdots$% 原稿“μ_1m_2g−”为写在 μ_2 前上方的插入小字

\noindent 都是减速（块推板、地阻板）% 以下原稿有一条贯穿页面的横线，其下为“③ 双速度”
%%%END
%%%BLOCK id=034-04 ch=03 sec=板块模型·有质板块·初速类 kind=method alt=01 cont=none y=0.68-0.75
\noindent ③ 双速度

%FIG p034_g 0.178 0.7005 0.35 0.772
\notefig[0.3]{figures/p034_g.png}
% 图内标注：物块(标“2”)在板上，向左箭头 V_0；板向右箭头 V_0；板上表面 μ_1，下表面 μ_2，板下标“1”；图下文字“底滑面糙”
\noindent 底滑面糙

\begin{align*}
a_2&=\mu_1g\\
a_1&=\frac{\mu_1m_2g+\mu_2(m_1+m_2)g}{m_1}
\end{align*}
\noindent $V_B=0$ 时\quad $t=\dfrac{V_0}{a_2}$\quad $S_B=\dfrac12V_0t$

\noindent $V_A=V_0-a_1t$

%FIG p034_h 0.664 0.6785 0.85 0.785
\notefig[0.33]{figures/p034_h.png}
% 图内标注：v-t 图，纵轴 V，正向 V_0、负向 -V_0，从 V_0 下降线标“1”，从 -V_0 上升线标“2”，交点后共同线标“1.2”，横轴 t
%%%END
%%%BLOCK id=034-05 ch=03 sec=板块模型·有质板块·外力类 kind=method alt=01 cont=none y=0.75-0.88
\textbf{外力类}

\noindent ① 恒力拉板

%FIG p034_i 0.05 0.797 0.475 0.876
\notefig[0.6]{figures/p034_i.png}
% 图内标注：长板(标“2”)上放物块(标“1”)，板受向右的 F；文字“恒力拉板”；v-t 图：纵轴 V、V_2、V_1，两条过原点的直线(上方的标 a_2，下方的标 a_1)，其间标“S_相”，虚线至 t_0，横轴 t
\[
X_{\text{相}}=\frac12a_2t_0^2-\frac12a_1t_0^2=L\quad\text{时分离}
\]
%%%END
%%%BLOCK id=034-06 ch=03 sec=板块模型·有质板块·外力类 kind=method alt=01 cont=none y=0.88-0.98
\noindent ② 恒力拉块

%FIG p034_j 0.049 0.884 0.49 0.978
\notefig[0.6]{figures/p034_j.png}
% 图内标注：文字“恒力拉块”；长板(标“2”)右端的物块受向左的 F；v-t 图：纵轴 V、V_1、V_2，两条过原点的直线(上方的标 a_1，下方的标 a_2)，其间标“S_相”，虚线至 t_0，横轴 t
\[
X_{\text{相}}=\frac12a_1t_0^2-\frac12a_2t_0^2=L\quad\text{时分离}
\]
%%%END
%%%BLOCK id=035-01 ch=03 sec=板块模型·力作用的最短时间 kind=method alt=01 cont=none y=0.06-0.24
\textbf{力作用的最短时间}

%FIG p035_a 0.085 0.152 0.23 0.20
\notefig[0.28]{figures/p035_a.png}
% 图内标注：长板(标“2”)上放物块(标“1”)，物块受向右的 F
%FIG p035_b 0.31 0.064 0.535 0.226
\notefig[0.42]{figures/p035_b.png}
% 图内标注：v-t 图(原稿在标题右侧)，纵轴 V，V_1、V_2，两条过原点的直线及转折，标 a、a_1、a_2，t_1、t_1+t_2，横轴 t；红笔：“S_相1”“S_相2”及两处红圈
\noindent $S_{\text{相}1}+S_{\text{相}2}=L$\quad $\dfrac12$ 水平\unclear{宽}$\times$铅垂高

\noindent $S_{\text{相}}=\dfrac12(V_2-V_1)(t_1+t_2)$
%%%END
%%%BLOCK id=035-02 ch=03 sec=板块模型·抽桌布 kind=method alt=01 cont=none y=0.23-0.45
\textbf{抽桌布}

%FIG p035_c 0.05 0.275 0.27 0.378
\notefig[0.4]{figures/p035_c.png}
% 图内标注：布(标“2”)上放物块 m(标“1”)，M；m 两侧标 L_1、L_2，接触面标 μ_1，F 向右；下方红笔画平行四边形(桌面)，内写 μ_2、“(m与桌)”
%FIG p035_d 0.295 0.245 0.51 0.368
\notefig[0.4]{figures/p035_d.png}
% 图内标注：v-t 图，纵轴 V、V_1、V_2，横轴 t、t_1、t_2；红字“板”“块”“S_1”“S_2”“S_3”，阴影；图上方文字“① 分离，② 1 停下”(“1”原稿像竖线)
\noindent 图注：① 分离\quad ② 1 停下

\begin{align*}
S_{\text{板}1}-S_{\text{块}1}&=S_1=L_1\\
S_2+S_3&=L_2
\end{align*}
\[
a_{\text{块}1}=+\mu_1g\qquad a_{\text{块}2}=-\mu_2g
\]
\[
F-\mu_1mg-f=Ma_{\text{板}}
\]
\noindent 若 $L_1=L_2$ 则\quad $\begin{aligned}S_2&=\dfrac{\mu_2}{\mu_1+\mu_2}\cdot\dfrac L2\\[4pt] S_3&=\dfrac{\mu_1}{\mu_1+\mu_2}\cdot\dfrac L2\end{aligned}$
%%%END
%%%BLOCK id=035-03 ch=03 sec=板块模型·共速与三块问题 kind=note alt=none cont=none y=0.45-0.49
\noindent 共速时 $f$ 消失\qquad 三块问题\qquad 板先和重物共速，形成整体后与轻块共速。
%%%END
%%%BLOCK id=035-04 ch=03 sec=板块模型·解题步骤 kind=method alt=none cont=none y=0.55-0.80
\textbf{步骤}

%FIG p035_e 0.085 0.598 0.22 0.65
\notefig[0.28]{figures/p035_e.png}
% 图内标注：板(标“2”)上放物块(标“1”)，物块受向右的 F，接触面 μ_1，地面 μ_2，板右端标 L
\[
\left\{\begin{array}{@{}l@{\quad}l@{\quad}l@{\quad}l@{\quad}l@{}}
\text{①} & \text{②} & \text{③} & \text{④} & \\
a_1= & V_1= & S_1= & V_{\text{相}}= & \text{恰好}\ d=S_1-S_2\ \ V_1=V_2\\
a_2=\ \cdots & V_2= & S_2= & a_{\text{相}}= & \\
\multicolumn{5}{@{}l}{S_{\text{相}}=S_1-S_2=L}\\
t=\ \cdots & & & &
\end{array}\right.
\]% ①②③④ 原稿为圈码，写在对应各列上方；“恰好 d=S_1-S_2  V_1=V_2”与第一行同高、位于右侧
%%%END
%%%BLOCK id=035-05 ch=03 sec=板块模型·抽桌布 kind=problem alt=01 cont=none y=0.80-1.00
\begin{problem}
抽桌布，\unclear{鱼}未掉求 $a$ 条件% “鱼”字不太确定(疑为鱼缸)；“a条件”写在下一行、与图重叠

%FIG p035_g 0.238 0.818 0.37 0.869
\notefig[0.3]{figures/p035_g.png}
% 图内标注：平行四边形(桌布)内标 μ_1、“m”小方块、“2”，向右箭头 a，下方标 μ_2
%FIG p035_f 0.12 0.8615 0.30 0.958
\notefig[0.36]{figures/p035_f.png}
% 图内标注：v-t 图，纵轴 V、V_2，旁边“a条件”；两条过原点的直线(较陡的与较缓的标 μ_1 g)，阴影区标 S_2；交点后下降线标 μ_2 g；S_1、S_1'(S_1 后带撇)，横轴 t、t'、t
\[
\frac{S_1}{S_1'}=\frac{\mu_2}{\mu_1}
\]
\begin{solution}
\noindent ① $a_1=\mu_1g$

\noindent ② $V_2=at$\qquad $V_1=a_1t$

\noindent\hspace*{2em}$S_2=\dfrac12at^2$\qquad $S_2=\dfrac12a_1t^2$% CHECK: 右式下标写得像 S_2，按上下文似应为 S_1，照原样转录

\noindent \hspace*{1em}布抽走后

\noindent ③ $a_1'=\mu_2g$

\noindent\hspace*{2em}$t'=\dfrac{V_1}{a_1'}$

\noindent\hspace*{2em}$S_1=V_1t'-\dfrac12a_1't'^2$% CHECK: 左端 S_1 原稿无撇(按语境似为 S_1')；“V_1t'”后有一小笔画

\[
\text{④}\ \left\{\begin{aligned}S_1+S_1'&=\frac L2\\ S_2-S_1&=\frac L2\end{aligned}\right.
\]
\[
S_1=\frac{\mu_2}{\mu_1+\mu_2}\cdot\frac L2\qquad S_1'=\frac{\mu_1}{\mu_1+\mu_2}\cdot\frac L2
\]
\end{solution}
\end{problem}
%%%END
%%%BLOCK id=036-01 ch=03 sec=板块模型·斜面叠放 kind=method alt=07 cont=prev y=0.00-0.115
\textbf{1.}
%FIG p036_a 0.07 0.0 0.26 0.075
\notefig[0.3]{figures/p036_a.png}

若底滑面糙 $(\mu_2\le\mu_1)$\quad 则共加共减速\quad 无相对位移\quad $a_1=a_2=g\sin\theta-\mu g\cos\theta$ % 右端被页边切断

若底糙面滑 $(\mu_2>\mu_1)$\quad 则物块会冲出去（跑得快）\quad 上快下慢，慢的会先停
%%%END
%%%BLOCK id=043-06 ch=03 sec=连接体·弹簧与分离 kind=problem alt=06 cont=none y=0.78-1.00
\begin{problem}[{[4]}]
%FIG p043_e 0.065 0.78 0.31 0.862
\notefig[0.3]{figures/p043_e.png}
% 图内标注：水平向右箭头 $F$ 推方块 $B$，$B$ 紧靠方块 $A$(两者上方各标 $m$)，$A$ 右侧连弹簧，弹簧另一端接右侧竖直墙；地面有斜线(粗糙)，弹簧下方标 $\mu$；图左上角有题号 [4]
$F$ 缓慢推动 $B$，压缩弹簧 $x_0$。此时 $A$、$B$ 静止，$B$ 不受摩擦。撤去 $F$ 后，$A$ 运动最大距离为 $4x_0$。求撤去 $F$ 瞬间，$A$、$B$ 的加速度。

\noindent $kx_0-\mu mg=2ma$% 原稿此处 k 写作大写 K

\noindent $a=\dfrac{kx_0}{2m}-\dfrac{\mu g}{2}$

\noindent $A$、$B$ 一起向左运动 $x$ 后\underline{分离}，求 $x$% “分离”二字下有红色下划线

\noindent 分离临界 $N=0$\quad $a_A=a_B=0$\qquad $k(x_0-x)=\mu mg$\qquad $\therefore x=x_0-\dfrac{\mu mg}{k}$
\end{problem}
%%%END
%%%BLOCK id=064-07 ch=03 sec=牛顿三定律 kind=summary alt=none cont=none y=0.69-0.79
\textbf{牛顿运动定律}
\begin{itemize}
\item 牛顿第一定律\quad 惯性 $\propto m$\quad 有相对性\quad 力是改变物体运动状态的原因 % 原稿“状”字为小字，插写在“态”字上方
\item 牛顿第二定律\quad $F=ma$\quad 单位制. 物理量和单位不同
\item 牛顿第三定律\quad 作用力、反作用力. \unclear{换元}. 同性反向等大
\end{itemize}
%%%END
%%%BLOCK id=064-08 ch=03 sec=牛顿三定律 kind=summary alt=none cont=none y=0.80-0.96
\textbf{对牛一、牛二、牛三定律的理解.}
\begin{itemize}
\item $F_{\text{合}}$与$v$同向 加速\quad 反向减速\quad 垂直圆运动 % CHECK: 末三字“垂直圆运动”，可能为“垂直：圆运动”
\item 惯性的相对性
  \begin{itemize}
  \item 铁球、水球、木球类.
  \end{itemize}
\item 矢量性\quad 瞬时性\quad 因果性 \quad （其上方有小字注：$a$与$F$）
\item 同一性\quad 同一物体、统一国际单位
\item 独立性\quad 每个力产生各自的加速度
\end{itemize}
%%%END
%%%BLOCK id=065-01 ch=03 sec=动力学两类基本问题 kind=summary alt=none cont=none y=0.04-0.15
\textbf{牛顿运动定律的基本运用}
\begin{itemize}
\item 动力学两类基本问题
  \begin{itemize}
  \item 已知受力情况求物体运动状态.
  \item 已知运动情况求物体受力情况\quad 以 $F_{\text{合}}=ma$ 为桥梁,
  \end{itemize}
\end{itemize}
%%%END
%%%BLOCK id=065-02 ch=03 sec=超重失重 kind=summary alt=none cont=none y=0.15-0.22
\begin{itemize}
\item 超重与失重\quad 向上超\ 下失
  \begin{itemize}
  \item 超加失减 $ma$.
  \item 太空人完全失重
  \end{itemize}
\end{itemize}
%%%END
%%%BLOCK id=065-03 ch=03 sec=瞬时加速度 kind=summary alt=none cont=none y=0.22-0.28
\textbf{瞬时加速度问题}
\begin{itemize}
\item 绳杆面、力\unclear{可}瞬间消失\quad 用牛二求 $a_{\text{瞬}}$
\item 弹簧、蹦床、橡皮\unclear{绳}、力保持
\end{itemize}
%%%END
%%%BLOCK id=065-04 ch=03 sec=超重失重 kind=note alt=none cont=none y=0.28-0.32
对超重与失重的理解.
%%%END
%%%BLOCK id=065-05 ch=03 sec=动力学两类基本问题 kind=note alt=none cont=none y=0.34-0.38
动力学的两类基本问题.
%%%END
%%%BLOCK id=065-06 ch=03 sec=动力学图像 kind=summary alt=none cont=none y=0.39-0.46
\textbf{动力学图像问题}
\begin{itemize}
\item $v$-$t$\quad $a$-$t$\quad $F$-$t$\quad $F$-$a$图.\quad 待定系数\ 或\ 利用数据
\end{itemize}
%%%END
%%%BLOCK id=065-07 ch=03 sec=斜面 kind=note alt=none cont=none y=0.47-0.50
三类光滑斜面上下滑时间比较.
%%%END
%%%BLOCK id=065-08 ch=03 sec=连接体与整体隔离 kind=summary alt=none cont=none y=0.52-0.66
\textbf{牛顿运动定律综合运用}
\begin{itemize}
\item 连接体问题
  \begin{itemize}
  \item 整体法.
  \item 隔离法
  \end{itemize}
\end{itemize}
%%%END
%%%BLOCK id=065-09 ch=03 sec=板块模型 kind=method alt=none cont=none y=0.67-0.78
\begin{itemize}
\item 滑块板块模型
  \begin{itemize}
  \item $F$、$a$、$v_{\text{共}}$、$t$、$x$、$Q$、快带慢、快后慢减 % 原稿“快带慢”与“快后慢减”之间上方有一小字 f（疑为摩擦力 f 的批注）
  \end{itemize}
\end{itemize}
%%%END
%%%BLOCK id=065-10 ch=03 sec=传送带 kind=method alt=none cont=none y=0.79-0.87
\begin{itemize}
\item 传送带模型
  \begin{itemize}
  \item 共速临界\quad 摩擦突变
  \end{itemize}
\end{itemize}
%%%END
%%%BLOCK id=065-11 ch=03 sec=临界极值 kind=summary alt=none cont=none y=0.89-1.00
\textbf{动力学中临界极值问题}
\begin{itemize}
\item 分离临界\quad $N=0$\quad $a_1=a_2$\quad $v_1=v_2$
\item 速度临界：合外力\unclear{为0}、收尾速度 % CHECK: “为0”二字连写形似“加”，按物理意义（收尾速度时合外力为零）读作“为0”
\item 滑动临界\quad 静摩擦达最大
\item 绳子松断临界：松 $T=0$；断 $T=\max$
\item \unclear{…} % 页底被照片边缘裁切，最后一行仅露出字的上端，无法辨认
\end{itemize}
%%%END
%%%BLOCK id=073-01 ch=03 sec=光滑斜面下滑时间 kind=derivation alt=01 cont=none y=0.05-0.30
% 页面最上方露出的封面（姓名栏）及右上角相邻页的倾斜小字「ρ=1.429 g·L^-1」属杂物/相邻页，未转录
\textbf{\unclear{光}滑面下滑时间} % CHECK: 首字被页边截断；原稿此处似无「斜」字（应为「光滑斜面下滑时间」）

\textbf{等高斜面}

%FIG p073_a 0.01 0.09 0.165 0.18
\notefig[0.25]{figures/p073_a.png}

\[
\left\{\begin{aligned}
L&=\tfrac{1}{2}at^{2}\\
a&=g\sin\theta\\
L&=\frac{h}{\sin\theta}
\end{aligned}\right.
\]
\[
t=\frac{1}{\sin\theta}\sqrt{\frac{2h}{g}}
\]
\[
t_1>t_2>t_3
\]
%%%END
%%%BLOCK id=073-02 ch=03 sec=光滑斜面下滑时间 kind=derivation alt=01 cont=none y=0.10-0.35
% 此处原稿无单独标题；图为底边长 d 相同、倾角不同的斜面（图中标 30°、两个 15° 增量），推测为「等底斜面」
%FIG p073_b 0.38 0.10 0.50 0.195
\notefig[0.25]{figures/p073_b.png}

\[
\left\{\begin{aligned}
L&=\tfrac{1}{2}at^{2}\\
a&=g\sin\theta\\
L&=\frac{d}{\cos\theta}
\end{aligned}\right.
\]
\[
t=\sqrt{\frac{4d}{g\sin 2\theta}}
\]

%FIG p073_c 0.625 0.24 0.80 0.35
\notefig[0.3]{figures/p073_c.png}

\[
t_1=t_3>t_2
\]
%%%END
%%%BLOCK id=073-03 ch=03 sec=等时圆 kind=method alt=01 cont=none y=0.34-0.68
\textbf{等时圆}

\textbf{同圆周同顶端的斜面}

%FIG p073_d 0.05 0.36 0.22 0.455
\notefig[0.28]{figures/p073_d.png}

\[
2R\sin\theta=\tfrac{1}{2}g\sin\theta\,t^{2}
\]
\[
t_1=t_2=t_3
\]

\textbf{同圆周同底端的斜面}

%FIG p073_e 0.04 0.455 0.21 0.555
\notefig[0.28]{figures/p073_e.png}

\[
t_1=t_2=t_3
\]

\textbf{斜面过公共切点且顶端在圆上}

%FIG p073_f 0.06 0.565 0.225 0.665
\notefig[0.28]{figures/p073_f.png}

\[
t_1=t_2=t_3
\]
%%%END
%%%BLOCK id=082-04 ch=03 sec=动力学图象·弹簧连接体 kind=method alt=01,06 cont=none y=0.64-0.80
%FIG p082_f 0.15 0.653 0.295 0.712
\notefig[0.28]{figures/p082_f.png}
\begin{align*}
a_A&=\frac{F-kx}{m}\ \downarrow\\
a_B&=\frac{kx}{m}\ \uparrow
\end{align*}
%FIG p082_g 0.485 0.632 0.685 0.734
\notefig[0.3]{figures/p082_g.png}
\begin{itemize}
\item $V_A=V_B$时\quad $a_A<a_B$
\item $a_A=a_B$时\quad $V_A>V_B$
\end{itemize}

$a_A=a_B$时\quad $\Delta V$最大

$V_A=V_B$时\quad $V_A\max$\quad $V_B\max$\quad $E_{p\text{弹}}\max$\quad $E_{\text{机}}\max$ % CHECK: V_A=V_B时“V_A max、V_B max”同时成立物理上存疑，照原样转录
%%%END
%%%BLOCK id=085-02 ch=03 sec=连接体·内力公式 kind=formula alt=none cont=none y=0.07-0.23
\textbf{内力公式}

\begin{align*}
\text{单一外力}\quad&T=\frac{m_1}{m_1+m_2}F\\[6pt]
\text{双反外力}\quad&T=\frac{m_1F_2+m_2F_1}{m_1+m_2}\\[6pt]
\text{双同外力}\quad&T=\frac{|m_1F_2-m_2F_1|}{m_1+m_2}
\end{align*}

同受的力不用写\quad 独受的力先合成后公式\quad 单体上合成\quad 力在不同物体不合成
%%%END
%%%BLOCK id=085-03 ch=03 sec=连接体·内力公式 kind=formula alt=none cont=none y=0.24-0.32
① % 原稿图中m_2下标注“μ=0”
%FIG p085_a 0.07 0.245 0.16 0.315
\notefig[0.25]{figures/p085_a.png}
\[
f_{\text{内}}=\frac{m_2F}{m_1+m_2}
\]
%%%END
%%%BLOCK id=085-04 ch=03 sec=连接体·内力公式 kind=formula alt=none cont=none y=0.24-0.32
④
%FIG p085_d 0.535 0.262 0.655 0.31
\notefig[0.25]{figures/p085_d.png}
\[
f_{\text{内}}=\frac{m_2F_1-m_1\bigl[F_2-\mu(m_1+m_2)g\bigr]}{m_1+m_2}
\]
%%%END
%%%BLOCK id=085-05 ch=03 sec=连接体·内力公式 kind=formula alt=none cont=none y=0.32-0.41
② % 原稿图中左侧标有f_地
%FIG p085_b 0.03 0.33 0.16 0.40
\notefig[0.25]{figures/p085_b.png}
\[
f_{\text{内}}=\frac{\mu(m_1+m_2)gm_1+Fm_2}{m_1+m_2}
\]
%%%END
%%%BLOCK id=085-06 ch=03 sec=连接体·内力公式 kind=formula alt=02 cont=none y=0.32-0.41
% 斜面上m_1叠在m_2上；图中标μ_1、μ_2；分子中“μ_2mgcosθ”被划去，其上方改写为μ_2(m_1+m_2)g(cosθ)
%FIG p085_e 0.505 0.325 0.61 0.395
\notefig[0.25]{figures/p085_e.png}
\[
f_{\text{内}}=\frac{m_1g\sin\theta\cdot m_2-\bigl[m_1g\sin\theta-\text{\struck{$\mu_2mg\cos\theta$}}\,\mu_2(m_1+m_2)g\cos\theta\bigr]m_1}{m_1+m_2}
\]
%%%END
%%%BLOCK id=085-07 ch=03 sec=连接体·内力公式 kind=formula alt=none cont=none y=0.41-0.49
③ % 原稿图中左侧标有f_地
%FIG p085_c 0.005 0.415 0.22 0.48
\notefig[0.3]{figures/p085_c.png}
\[
f_{\text{内}}=\frac{m_1\bigl[F-\mu(m_1+m_2)g\bigr]}{m_1+m_2}
\]
%%%END
%%%BLOCK id=085-08 ch=03 sec=连接体·内力公式 kind=formula alt=none cont=none y=0.41-0.52
⑤
%FIG p085_f 0.565 0.414 0.68 0.462
\notefig[0.25]{figures/p085_f.png}
\[
f_{\text{内}}=\frac{m_1\mu_2(m_1+m_2)g}{m_1+m_2}
\]
底滑面粗糙

应$<\mu_1m_1g$

故$\mu_1>\mu_2$\ \unclear{共}减速
%%%END
%%%BLOCK id=085-09 ch=03 sec=连接体·内力公式 kind=formula alt=none cont=none y=0.53-0.67
⑥
%FIG p085_g 0.04 0.535 0.19 0.61
\notefig[0.25]{figures/p085_g.png}
\begin{align*}
f_{12}&=\frac{\bigl[F_1-\mu(m_1+m_2+m_3)g\bigr](m_2+m_3)-(F_3-F_2)m_1}{m_1+m_2+m_3}\\[8pt]
f_{23}&=\frac{\bigl[F_1-F_2-\mu(m_1+m_2+m_3)g\bigr]m_3-F_3(m_1+m_2)}{m_1+m_2+m_3}
\end{align*}
%%%END
%%%BLOCK id=085-10 ch=03 sec=连接体·内力公式 kind=formula alt=none cont=none y=0.67-0.82
⑦ % 原稿图中m_2、m_1被一小椭圈圈起，整个系统又被一大椭圈圈起
%FIG p085_h 0.06 0.668 0.275 0.755
\notefig[0.35]{figures/p085_h.png}
\[
T=\frac{(m_1+m_2)F}{m_1+m_2+m_3}
\]
\[
f_{12}=\frac{m_2\bigl[F-\mu(m_1+m_2+m_3)g\bigr]}{m_1+m_2+m_3}
\]
%%%END
%%%BLOCK id=085-11 ch=03 sec=连接体·内力公式 kind=formula alt=none cont=none y=0.82-1.00
⑧
%FIG p085_i 0.05 0.805 0.30 0.884
\notefig[0.35]{figures/p085_i.png}
\[
f_{\text{左}}=\frac{2mF}{6m}\qquad\qquad f_{\text{右}}=\frac{4m}{6m}F
\]
对$1\,2\,3$\quad $T=\dfrac{3m\cdot\mu mg}{3m+m}$

改\quad $f_{\text{右}}>2f_{\text{左}}$\quad 左先滑\quad $F$到$3f_{\text{左}}$时\quad $C,D$相对\unclear{滑} % 行末被页面右缘截断

$f_{\text{右}}<2f_{\text{左}}$\quad 右先滑\quad $F$到$\dfrac32f_{\text{右}}$时\quad $A,B$相对\unclear{滑} % 行末被页面右缘截断
%%%END
%%%BLOCK id=086-01 ch=03 sec=板块模型·滑动临界 kind=method alt=none cont=none y=0.03-0.34
\textbf{滑动临界}

%FIG p086_a 0.085 0.055 0.21 0.11
\notefig[0.25]{figures/p086_a.png}
\[ f=\frac{m_2F+\mu_2(m_1+m_2)g\cdot m_1}{m_1+m_2}<\mu_1m_1g\qquad\text{为共加速条件} \]

%FIG p086_b 0.09 0.125 0.21 0.19
\notefig[0.25]{figures/p086_b.png}
\[ f=\frac{\mu_2(m_1+m_2)g\cdot m_1}{m_1+m_2}<\mu_1m_1g\qquad\text{为共加速条件} \]

%FIG p086_c 0.11 0.2 0.215 0.25
\notefig[0.25]{figures/p086_c.png}

\noindent 若 $\mu_1m_1g<\mu_2(m_1+m_2)g$ 则2不动 \\
若 $\mu_1m_2g>\mu_2(m_1+m_2)g$ 则2也动 % CHECK: 此行 m 的下标手写看起来是2（物理上应与上一行一样是 m_1）
\[ f=\frac{Fm_2+\mu_2(m_1+m_2)g\,m_1}{m_1+m_2}\ge\mu_1m_1g\ \text{时}\quad\text{相对滑动} \]
%%%END
%%%BLOCK id=086-03 ch=03 sec=连接体·整体隔离 kind=problem alt=none cont=none y=0.45-0.745
\begin{problem}[2、]
火车从东拉以 $a$，PQ 间拉力为 $F$；从西拉以 $\frac23a$，PQ 间拉力为 $F$，求节数
\begin{solution}
%FIG p086_e 0.115 0.512 0.33 0.59
\notefig[0.3]{figures/p086_e.png}
\[ n=x+y\ \text{节} \]
\[ F=xm_0a \]
% m 后的小圆圈为下标 0（每节车厢质量 m_0）
%FIG p086_f 0.055 0.605 0.28 0.665
\notefig[0.3]{figures/p086_f.png}
\[ F=\frac23ay\,m_0 \]
\[ \begin{cases} x=\dfrac23y\\ x+y=n\end{cases}\qquad n=\frac53y \]
\[ 3n=5y\qquad \frac n5=\frac y3 \]
\[ \begin{aligned} n&\text{ 为 }5k_1,\\ y&\text{ 为 }3k_1,\\ x&\text{ 为 }2k_1\end{aligned} \]
\end{solution}
\end{problem}

\noindent 动车拖车相间，任意拖动\,间内力为0，任意动拖间内力为 $\dfrac F2$ \\
力为$F$\ 力为0
% CHECK: “任意拖动”疑为“任意拖拖”；第二行行首“力为F 力为0”是附注，行首字被页边遮挡，含义不明（疑为“动动间内力为F，拖拖间内力为0”）
%%%END
%%%BLOCK id=086-04 ch=03 sec=连接体·整体隔离 kind=method alt=none cont=none y=0.745-0.84
\noindent 3、整体/隔离法。几个物体几组独立方程，整体列1组，隔离\unclear{物体}少的列$n-1$组，横竖列1个，倾斜列\unclear{2}组 \quad $\begin{cases}F_x=\cdots\\ F_y=\cdots\end{cases}$
% “n-1”写在行上方，行内“列”与“组”之间有涂改/划线痕迹；“倾斜列2组”末尾与右侧大括号被页边截断

\medskip
\noindent 关联加速度\quad 一根绳上 $a$ 相\,等（方向可不共线）\qquad 有对称少1组，有临界临界 \\
动滑轮加速度、速度受力有倍数关系 \\
(\underline{不用}整体法，滑轮也得分析)
% “相”与“等”之间有涂改液；“有临界临界”行末被页边截断，后面可能还有字；“动滑轮…”一行是写在行间的小字
%%%END
%%%BLOCK id=086-05 ch=03 sec=连接体·关联加速度 kind=problem alt=none cont=none y=0.84-0.99
\begin{problem}
%FIG p086_g 0.058 0.848 0.15 0.97
\notefig[0.25]{figures/p086_g.png}
\begin{solution}
\[ \begin{cases} T-mg=ma\\ T-Mg+N=0\end{cases} \]
\[ N=Mg-ma-mg\qquad\text{系统失重} \]
\end{solution}
\end{problem}
%%%END
%%%BLOCK id=086-06 ch=03 sec=连接体·整体隔离 kind=problem alt=02 cont=none y=0.83-0.99
\begin{problem}
%FIG p086_h 0.455 0.858 0.61 0.952
\notefig[0.3]{figures/p086_h.png}
BC 静止，求 $a_A$
\begin{solution}
对A
\[ \begin{cases} mg\sin30^\circ-f=ma\\ 1.5mg\sin30^\circ=T+f\\ T=mg\end{cases} \]
% CHECK: 第二式按受力应为 1.5mg sin30°+f=T 才能得到 a_A=g/4；此处照抄原稿 “=T+f”
$\therefore a_A=\dfrac g4$
\end{solution}
\end{problem}
%%%END
%%%BLOCK id=087-01 ch=03 sec=超重失重·系统牛二 kind=method alt=none cont=none y=0.06-0.20
\textbf{流动问题}\quad 谁$\rho$小谁朝运动方向走。
% CHECK: “ρ”手写形似大写 P，按物理应为密度 ρ
%FIG p087_a 0.055 0.095 0.21 0.15
\notefig[0.25]{figures/p087_a.png}
\noindent 若为木球/气泡/$\cdots$\quad 球往右 \\
若为铁球\quad 球往左

%FIG p087_b 0.61 0.12 0.70 0.195
\notefig[0.25]{figures/p087_b.png}
\noindent 木球上浮\quad $N$小，失重(水)向下加\unclear{速}
% 此行行末被页边截断；括号形状不确定
\\
铁球下沉时 $N$ 小，铁失重
%%%END
%%%BLOCK id=087-02 ch=03 sec=系统牛顿第二定律 kind=formula alt=none cont=none y=0.20-0.27
\noindent 系统牛顿第二定律、不分析内力\quad 安培力\quad 不是内力是磁场施力

\noindent 非共加速\quad $F_{\text{合}}=m_1a_1+m_2a_2+\cdots$\quad 整体合外力提供部分加速度 \\
共加速\quad $F_{\text{合}}=(m_1+\cdots+m_n)a$
%%%END
%%%BLOCK id=087-03 ch=03 sec=超重失重 kind=formula alt=none cont=none y=0.27-0.345
\[ a\text{ 向}\left\{\begin{array}{ll}\text{上 超}&\quad N=mg+ma\\ \text{下 失}&\quad N=mg-ma\end{array}\right.\qquad \begin{array}{l}\text{上 超 下 失}\\ \text{超 加 失 减}\end{array} \]
%%%END
%%%BLOCK id=087-04 ch=03 sec=超重失重·系统牛二 kind=problem alt=none cont=none y=0.35-0.47
\begin{problem}
%FIG p087_c 0.07 0.35 0.145 0.44
\notefig[0.25]{figures/p087_c.png}
环上升至 $l$ 时 $N_{\text{地}}$
% “N地”的“地”字有重描痕迹
\begin{solution}
\[ (M+m)g-N=ma \]
\[ v_0^2=2al \]
\end{solution}
\end{problem}
%%%END
%%%BLOCK id=087-05 ch=03 sec=系统牛顿第二定律 kind=problem alt=none cont=none y=0.485-0.575
\begin{problem}
%FIG p087_d 0.10 0.485 0.31 0.545
\notefig[0.3]{figures/p087_d.png}
\begin{solution}
\[ f_{\text{地}}=ma\cos\theta \]
\[ (M+m)g-N_{\text{地}}=ma\sin\theta \]
\end{solution}
\end{problem}
%%%END
%%%BLOCK id=087-06 ch=03 sec=系统牛顿第二定律 kind=problem alt=none cont=none y=0.57-0.65
\begin{problem}
%FIG p087_e 0.145 0.57 0.35 0.64
\notefig[0.3]{figures/p087_e.png}
\begin{solution}
\[ \begin{cases} f=m_2a\cos\theta\\ m_2g\sin\theta-m_3g=(m_2+m_3)a\\ (m_1+m_2+m_3)g-N=m_1a\cos\theta-m_3a\end{cases} \]
% CHECK: 第三式右边手写为 m_1 a cosθ − m_3 a（按受力分析 m_2 的竖直分加速度应为 m_2 a sinθ），照抄原稿
\end{solution}
\end{problem}
%%%END
%%%BLOCK id=087-07 ch=03 sec=系统牛顿第二定律 kind=problem alt=none cont=none y=0.65-0.71
\begin{problem}
$a=g\sin\theta$
%FIG p087_f 0.18 0.648 0.30 0.705
\notefig[0.25]{figures/p087_f.png}
$\theta$ 可变，求 $f_{\text{地}\,\max}$
\begin{solution}
\[ f=ma\cos\theta=mg\cdot\frac12\sin2\theta_{\max}=\frac{mg}{2}\qquad \theta=45^\circ\text{ 时取等。} \]
% “max”写在 sin2θ 右下角，意为 f_max
\end{solution}
\end{problem}
%%%END
%%%BLOCK id=087-09 ch=03 sec=连接体·整体隔离 kind=method alt=none cont=none y=0.90-0.94
\noindent 能整体先整体，剩余物体 $n-1$\quad 不能整体全隔离，关联加速未知力
%%%END
%%%BLOCK id=089-01 ch=03 sec=非惯性系·惯性力 kind=knowledge alt=none cont=none y=0.04-0.125
\noindent $\to a$\quad $\leftarrow ma$
% 位于标题上方的小字

\noindent 非惯性系。反向施\unclear{加}$ma$ 使之平衡（假想）\quad 施加惯性力\ 并不存在
% “非惯性系”原稿被菱形框住；“施加”的“加”字形似“向”
%%%END
%%%BLOCK id=114-04 ch=03 sec=连接体与整体隔离 kind=problem alt=90 cont=none y=0.58-0.72
\begin{problem}[4]
\textbf{$\sin2\theta=2\sin\theta\cos\theta$ 及辅助角公式}

% 图p114_f内含标注：g sinθ、μ=0（或M=0，原稿字形难辨）、m、M、θ
%FIG p114_f 0.025 0.645 0.24 0.71
\notefig[0.3]{figures/p114_f.png}

当 $\theta=?$ 时 $f_{\text{地}}$ max\qquad 系统牛二律
\begin{solution}
\[f_{\text{地}}=M\cdot0+ma\cos\theta=mg\sin\theta\cos\theta\qquad \theta=45^\circ\ \text{时取极大值}\]
\end{solution}
\end{problem}
%%%END
%%%BLOCK id=114-05 ch=03 sec=斜面 kind=problem alt=02,90 cont=none y=0.72-0.90
% 原稿序号处“4”被涂改为“5”
\begin{problem}[5]
% 图p114_g内含标注：N、F、f、mg、φ、30°、μ=tan30°
%FIG p114_g 0.085 0.735 0.215 0.84
\notefig[0.3]{figures/p114_g.png}

$\to a=\dfrac{g}{2}$\quad 求最小拉力 $F$
\begin{solution}
\[\left\{\begin{array}{l}F\cos\alpha-mg\sin30^\circ-f=ma\\ N+F\sin\alpha=mg\cos30^\circ\\ f=\mu N\end{array}\right.\qquad F=\frac{mg\sin30^\circ+\mu mg\cos30^\circ+ma}{\sqrt{1+\mu^2}\,\sin(\alpha+\varphi)}\]
\[\tan\varphi=\frac{1}{\mu}\]
当 $\sin(\alpha+\varphi)=1$ 时 $F$ 最小
\end{solution}
\end{problem}
%%%END
%%%BLOCK id=119-01 ch=03 sec=弹簧·物块落在弹簧上 kind=knowledge alt=06 cont=none y=0.07-0.21
% 标题“物块”二字前有一个难辨的小符号，未转录
\textbf{\underline{物块落在弹簧上}}

%FIG p119_a 0.075 0.10 0.145 0.188
\notefig[0.25]{figures/p119_a.png}
% 图中标注：弹簧上方一个圆圈(物块)，弹簧旁自上而下标 A、O、B 三个位置

\begin{itemize}
\item[A] $kx<mg$\quad $v\uparrow$\quad $a$向下$\downarrow$
\item[O] $kx=mg$\quad $v_{\max}$
\item[B] $kx>mg$\quad \struck{\unclear{…}}\ $v\uparrow$\quad $a$向上$\downarrow$ % CHECK: B 行 v、a 箭头方向与物理预期相反(预期 v 减小、a 向上增大)；kx>mg 后另有一处被涂改的字
\end{itemize}

$v_{\max}$不在刚碰处

%FIG p119_b 0.405 0.121 0.592 0.198
\notefig[0.25]{figures/p119_b.png}
% 图中标注：纵轴 v，$V_m$ 等标注，曲线上一点 D，横轴 t

\[
\left\{\begin{aligned}
&\text{有}h\text{释放}\quad a_{\text{下}}>g\\
&\text{无}h\text{释放}\quad a_{\text{下}}=g\\
&\text{提着释放}\quad a_{\text{下}}<g
\end{aligned}\right.
\qquad\text{动能定理证明}
\]
%%%END
%%%BLOCK id=119-02 ch=03 sec=动力学临界问题 kind=knowledge alt=none cont=none y=0.21-0.28
变速临界\quad $a=0$\ $v_{\max}$\quad $v=0$\ $a_{\max}$

分离临界\quad $N=0$\ $a_1=a_2$\ $v_1=v_2$\quad 不分离\ $N\ge 0$

滑动临界、整体隔离法
%%%END
%%%BLOCK id=164-05 ch=03 sec=动力学图像·外力F作用 kind=note alt=02 cont=none y=0.79-0.86
% “在水平面上”下方原稿为波浪下划线
外力$F$作用$\cdots$\ \underline{在水平面上}\ %
%FIG p164_f 0.39 0.795 0.51 0.85
\notefig[0.2]{figures/p164_f.png}
只能算出$F_{\text{水平}}$\quad 无法算出$F$(未知方向)
%%%END
%%%BLOCK id=199-07 ch=03 sec=有空气阻力的竖直上抛 kind=formula alt=01 cont=none y=0.86-0.95
有空气阻力

\notefig[0.1]{figures/p199_f.png}
% 图为两个带小圆圈（物体）的竖直箭头：上方箭头向上，下方箭头向下

\[
a_1=\frac{f+mg}{m}
\]
\[
a_2=\frac{mg-f}{m}\ \,{}^{a\text{不同}}
\]
%%%END
